% BORRADOR (comentarios del director, 19-oct-2026): reemplaza MainMatter/Cap1/M02-Chapter1.tex. Borrar esta línea al pasarlo.

\section{Planteamiento del problema}
\label{sec:problema}

La búsqueda de empleo remoto se reparte entre muchas plataformas: agregadores generalistas, portales especializados y bolsas de empleo de las propias empresas. Cada una tiene su catálogo y sus filtros, y muchas resuelven la búsqueda de la misma manera: comparan los términos de la consulta con los del título o la descripción de la oferta. Ese mecanismo funciona cuando quien busca y quien publica nombran el cargo con las mismas palabras, y falla cuando no lo hacen.

La literatura de emparejamiento de empleo documenta esta limitación desde hace tiempo. \citet{HaThuc2017} describen cómo LinkedIn pasó de una búsqueda de candidatos por palabras clave a una búsqueda por ejemplos, porque las palabras clave no capturaban lo que el reclutador buscaba. \citet{Bocharova2023} atribuyen a la variación léxica entre títulos equivalentes, como \enquote{Software Engineer} y \enquote{Backend Developer}, una pérdida sistemática de coincidencias, y proponen representaciones semánticas entrenadas para el dominio. La revisión bibliométrica de \citet{RojasGaleano2022} muestra que esta limitación es uno de los motivos centrales de la investigación en inteligencia artificial aplicada al emparejamiento entre ofertas y currículos.

Las grandes plataformas ya adoptaron alternativas. LinkedIn recupera ofertas con representaciones densas aprendidas \citep{Shen2024} y usa un modelo de lenguaje para interpretar la consulta y extraer de ella atributos estructurados \citep{Liu2025}. Otros trabajos usan modelos de lenguaje para generar recomendaciones de empleo \citep{ZhengGIRL2023}. Estas soluciones, sin embargo, se entrenan y se evalúan con datos de interacción propietarios, a los que no tiene acceso un agregador pequeño.

La generación aumentada por recuperación (\textit{Retrieval-Augmented Generation}, RAG) \citep{Lewis2020} ofrece una alternativa construida con componentes disponibles. Un sistema RAG recupera los fragmentos más cercanos a la consulta en un espacio de \textit{embeddings} y le pide a un modelo de lenguaje que redacte la respuesta a partir de ellos. La revisión de \citet{Gao2023} distingue un RAG \emph{ingenuo}, con una sola búsqueda densa, de un RAG \emph{avanzado}, que agrega optimizaciones como el reordenamiento de los resultados, y discute además la recuperación híbrida, que combina representaciones léxicas y densas. Cada etapa promete mejorar la recuperación, y cada una tiene un costo en latencia y en complejidad.

Dos vacíos motivan este trabajo. El primero es de evidencia. Para la búsqueda de empleo no se sabe cuánto mejora un pipeline RAG la recuperación frente a una línea base léxica fuerte como BM25, ni qué componente explica esa mejora. La pregunta no es trivial: en la comparación de \citet{Thakur2021} sobre 18 colecciones, BM25 resultó una línea base robusta que varios modelos densos no superaron fuera de su dominio de entrenamiento, y los modelos de reordenamiento obtuvieron los mejores resultados a un alto costo computacional. El segundo vacío es de método. Un sistema RAG responde con fluidez aunque su recuperación, sus datos o su enrutamiento fallen, porque el modelo de lenguaje redacta una respuesta plausible con lo que recibe. \citet{Barnett2024} identifican siete puntos de falla en la ingeniería de estos sistemas y concluyen que su validación solo es posible durante la operación. Falta, entonces, un protocolo que revele esas fallas antes, componente por componente.

Este trabajo estudia ambos vacíos con un caso concreto: JobAI, una plataforma de búsqueda de empleo remoto que agrega ofertas de tres fuentes públicas y responde consultas en lenguaje natural mediante un pipeline RAG. JobAI es un artefacto preexistente, desarrollado por el autor sin un protocolo de evaluación formal. Sirve como caso de estudio porque su pipeline admite activar y desactivar componentes de forma aislada.

\section{Pregunta de investigación}

\begin{quote}
\textit{¿En qué medida un pipeline de Retrieval-Augmented Generation mejora la recuperación de ofertas de empleo frente a la búsqueda por palabras clave, qué componentes del pipeline explican esa mejora y con qué costo en latencia?}
\end{quote}

La forma de evaluar un artefacto que ya existía es una decisión metodológica, no una pregunta de investigación. Se trata en el Capítulo~\ref{cap:metodologia}.

\section{Objetivos}

\subsection{Objetivo general}

Evaluar en qué medida un pipeline de \textit{Retrieval-Augmented Generation} aplicado a la búsqueda de empleo remoto mejora la calidad de la recuperación frente a la búsqueda por palabras clave, e identificar el aporte y el costo de cada uno de sus componentes, usando la plataforma JobAI como caso de estudio.

\subsection{Objetivos específicos}

\begin{enumerate}
    \item Caracterizar, a partir de la literatura, las limitaciones de la búsqueda léxica en portales de empleo y los componentes de RAG que se proponen para superarlas: recuperación densa, filtros, búsqueda híbrida, reordenamiento y modelos multilingües.
    \item Diseñar un protocolo de evaluación reproducible: corpus congelado, consultas por categoría, juicios de relevancia con un modelo de lenguaje como juez validado contra juicios humanos, métricas de recuperación, fidelidad y latencia, y análisis estadístico.
    \item Aplicar el protocolo para mejorar el pipeline de JobAI en dos iteraciones de diseño y evaluación: (a) corrección de las fallas que el protocolo detectó y (b) incorporación de componentes de recuperación.
    \item Cuantificar el aporte de cada componente frente a BM25 y a la búsqueda de la plataforma, en calidad de la recuperación y fidelidad de la respuesta, junto con su costo en latencia y el efecto del motor de almacenamiento vectorial.
\end{enumerate}

\section{Justificación}

\paragraph{Práctica.} Quien busca empleo pierde tiempo revisando ofertas irrelevantes y no ve ofertas compatibles que el buscador no reconoce. Muchas plataformas anuncian hoy \enquote{búsqueda con inteligencia artificial} sin medir qué ganan con ella. Una evaluación por componentes indica cuáles vale la pena implementar y cuánto cuestan, una información útil para cualquier plataforma pequeña que considere un pipeline RAG.

\paragraph{Teórica.} La investigación en emparejamiento de empleo se ha concentrado en modelos de \textit{ranking} y recomendación entrenados con datos de interacción \citep{HaThuc2017,Shen2024,Bocharova2023}. La literatura de RAG, por su parte, se evalúa sobre todo en colecciones de preguntas y respuestas de dominio general \citep{Gao2023}. Este trabajo cruza ambas líneas con un caso evaluado de forma reproducible: mide el aporte de la búsqueda densa, los filtros, la búsqueda híbrida, el reordenamiento y un modelo multilingüe sobre ofertas reales, frente a una línea base léxica fuerte.

\paragraph{Metodológica.} El protocolo propuesto separa las decisiones de configuración del conjunto de prueba, valida al juez automático contra juicios humanos y registra cada respuesta, juicio y versión del código. Además, al aplicarse a un artefacto existente, permite observar qué fallas de un sistema RAG solo se hacen visibles con una evaluación sistemática. Esa caracterización es útil para quien evalúe un sistema RAG en otro dominio.

\section{Alcance y limitaciones}
\label{sec:alcance}

El trabajo evalúa la configuración de referencia de JobAI (Capítulo~\ref{cap:artefacto}) sobre un corpus congelado de ofertas de empleo remoto de tres fuentes públicas: Remotive, Jobicy y Findwork. Mide la calidad de la recuperación, la fidelidad de las respuestas, la latencia por etapa y el comportamiento de dos motores de almacenamiento vectorial, según el protocolo del Capítulo~\ref{cap:metodologia}.

Quedan fuera del alcance:
\begin{itemize}
    \item La evaluación con usuarios reales: el trabajo mide la calidad técnica del pipeline, no la experiencia de uso ni la adopción.
    \item El diálogo de varios turnos. La plataforma incluye el historial reciente de la conversación en la generación, pero la recuperación usa solo la consulta actual, y la evaluación trata cada consulta de forma aislada.
    \item La personalización por currículo, que se desactiva durante la evaluación para no confundir su efecto con el de la recuperación.
    \item La generalización a otros dominios, idiomas de las ofertas o modelos de \textit{embeddings} de otras familias.
    \item La comparación con todos los motores vectoriales del mercado: se comparan pgvector y Qdrant como representantes de una extensión sobre un motor relacional y de un motor dedicado.
    \item La calidad de los datos de origen, que se toman como dados, aunque el trabajo documenta los defectos que afectan a la evaluación.
    \item El despliegue público de la plataforma, cuya configuración difiere de la evaluada.
\end{itemize}

\section{Marco teórico conceptual}
\label{sec:marco-teorico}

Esta sección introduce el vocabulario necesario para seguir el documento. El Capítulo~\ref{cap:estado-arte} revisa la literatura de cada concepto y su relación con JobAI.

\subsection{Recuperación léxica y desajuste de vocabulario}

Un sistema de recuperación léxica representa cada documento por sus términos y lo puntúa según los términos que comparte con la consulta. El modelo de referencia es Okapi BM25 \citep{Robertson2009}, que pondera cada término coincidente por su rareza en la colección y normaliza por la longitud del documento. Su limitación es el \emph{desajuste de vocabulario}: dos textos equivalentes que no comparten términos no se relacionan. En el reclutamiento aparece como sinonimia entre habilidades (\enquote{JS} y \enquote{JavaScript}), variación entre títulos de cargo y diferencias entre cómo una persona describe lo que busca y cómo una oferta describe el cargo. La búsqueda por subcadena, del tipo \texttt{ILIKE} de SQL, es una forma aún más simple de búsqueda léxica: exige que el texto de la consulta aparezca literalmente.

\subsection{Modelos de lenguaje de gran escala}

Los modelos de lenguaje de gran escala (\textit{Large Language Models}, LLM) se preentrenan sobre grandes volúmenes de texto y producen lenguaje fluido a partir de instrucciones. Su conocimiento queda fijado en sus parámetros al terminar el entrenamiento. No conocen datos posteriores ni datos privados, como el catálogo de una plataforma de empleo, y tienden a \enquote{alucinar}, es decir, a producir afirmaciones plausibles pero no sustentadas cuando les falta información \citep{Gao2023}.

\subsection{Retrieval-Augmented Generation}

RAG \citep{Lewis2020} combina un componente de recuperación, que busca en una base de conocimiento externa los fragmentos más relevantes para la consulta, con un LLM que redacta la respuesta a partir de ellos. Así la respuesta se apoya en información verificable y actualizable. \citet{Gao2023} distinguen tres niveles:
\begin{itemize}
    \item \emph{RAG ingenuo}: indexación, una única búsqueda y generación.
    \item \emph{RAG avanzado}: agrega optimizaciones antes y después de la recuperación, como la reformulación de la consulta y el reordenamiento de los resultados.
    \item \emph{RAG modular}: organiza el sistema en componentes intercambiables, con flujos condicionales como el enrutamiento entre estrategias.
\end{itemize}
La \emph{búsqueda híbrida}, que la misma revisión trata entre los recuperadores, combina una lista léxica y una densa; un método habitual es la fusión por rangos recíprocos (RRF), que suma, para cada documento, el inverso de su posición en cada lista \citep{Cormack2009}. El \emph{reordenamiento} aplica un modelo más costoso, típicamente un \textit{cross-encoder} que lee juntas la consulta y cada candidata, a las primeras posiciones de esa lista \citep{Nogueira2019}.

Evaluar un sistema RAG exige medir por separado la recuperación y la generación. Para la segunda, \citet{Es2024} proponen métricas sin respuesta de referencia, como la \emph{fidelidad}, que verifica si cada afirmación de la respuesta se sostiene en el contexto recuperado, y la \emph{relevancia de la respuesta}. El Capítulo~\ref{cap:metodologia} retoma estas dos.

\subsection{Bases de datos vectoriales y búsqueda por similitud}

Un \textit{embedding} es una representación vectorial de un texto en un espacio donde la cercanía entre vectores, típicamente medida con la similitud coseno, refleja la cercanía semántica entre los textos. Recuperar los fragmentos más relevantes equivale a buscar los vecinos más cercanos del vector de la consulta. Como la búsqueda exacta es costosa a gran escala, se usan algoritmos de búsqueda \emph{aproximada} de vecinos (ANN). El más influyente es HNSW \citep{Malkov2020}, un grafo jerárquico de navegación con búsquedas en tiempo aproximadamente logarítmico. Estos algoritmos se ofrecen como extensión de un motor relacional, como pgvector sobre PostgreSQL, o dentro de un motor dedicado a vectores, como Qdrant o Milvus. Cuál conviene depende de la escala y del uso de filtros; el Capítulo~\ref{cap:estado-arte} revisa la evidencia y el Capítulo~\ref{cap:resultados} la contrasta con datos reales y sintéticos.

\subsection{Design Science Research}
\label{sec:marco-teorico-dsr}

Design Science Research (DSR) es un paradigma de investigación orientado a crear y evaluar con rigor artefactos que resuelven problemas relevantes: constructos, modelos, métodos o instanciaciones de software. \citet{Hevner2004} lo formalizan en siete lineamientos, entre ellos que el artefacto debe evaluarse con rigor y que la investigación debe producir una contribución verificable. \citet{Peffers2007} lo operacionalizan en seis actividades: identificación del problema, definición de objetivos, diseño y desarrollo, demostración, evaluación y comunicación. El proceso admite distintos puntos de entrada, entre ellos uno centrado en el diseño y desarrollo, que parte de un artefacto existente. El Capítulo~\ref{cap:metodologia} describe su aplicación a JobAI.
